% ============================================================
% kbs_cost.tex — KBS journal: Acquisition cost (Section 7).
% Section label: sec:cost. RQ4 (plan). The E4 instrument is built and
% validated (phase1/scripts/export_cost_study.py, score_cost_study.py);
% results pending the cataloger panel -> marked \TODO{E4}.
% Bibliography: kbs_ke_references.bib + paper/references.bib.
% ============================================================

\section{Cost of Machine- vs.\ Expert-Acquired Knowledge}
\label{sec:cost}

Validity (Section~\ref{sec:valid}) and utility (Section~\ref{sec:util})
establish that machine-acquired subject knowledge is useful. Whether a
knowledge manager should deploy it depends on cost: expert time is the
scarce resource in knowledge acquisition \cite{feigenbaum_knowledge_engineering},
and the relevant comparison is not machine-versus-nothing but
machine-versus-expert for the same knowledge. This section reports the
design and protocol; results are pending the cataloger panel
\TODO{E4: report panel results}.

\subsection{Design}

We measure acquisition cost in a counterbalanced panel study
\TODO{E4}, following the classical framing that knowledge acquisition cost
is expert time \cite{studer_knowledge_engineering}. Each expert cataloger
processes 80 records in two conditions: \emph{manual} acquisition---assign
subject headings and a classification from a sparse record (title, author,
year), the same input the acquisition component receives---and
\emph{verification} acquisition---review and edit a machine-proposed
heading set and classification for a record of the same type. Conditions are
disjoint within a cataloger and counterbalanced across catalogers; neither
condition reveals the professional gold. For each record we instrument
elapsed time and record the final heading set.

Two comparisons come out of this design. \emph{Time}: mean seconds per
record under manual versus verification acquisition, per cataloger and
pooled. \emph{Quality}: the final heading sets from both conditions scored
against validated gold with the multi-level rubric of
Section~\ref{sec:valid}, so we can ask whether verification saves time at
no quality cost, or quantify the trade. A secondary measure is the
\emph{edit rate} in the verification condition---how often the expert
changed the machine draft---which separates ``the draft was good and the
expert kept it'' from ``the expert repaired a bad draft'' (the latter being
hidden cost). The instrument is released and fully scripted
\TODO{E4: report results; attach the released instrument}.

\subsection{What the design controls}

The comparison is fair by construction on the input side: both conditions
see the same sparse record, so any time difference is attributable to the
draft, not to different information. The comparison is deliberately
conservative on the machine side: the acquisition component is unassisted
(single model, sparse-record conditioning; multi-model and
full-text-conditioned variants are reported in Section~\ref{sec:valid}
\TODO{E3}), so the measured verification time is an upper bound on the
time an assisted workflow would need. The panel is small (two to three
catalogers), so per-cataloger results are reported rather than pooled only,
and the cost function is treated in Section~\ref{sec:synth} as a pilot
estimate whose sensitivity is explored rather than over-claimed.
